\scanpage{13}
la connaissance des $\mu_i$. Si on se donne les \inserted{entiers} $\nu_j\geq\unsure{0}$ \deleted{$\unsure{\mathrm{rg}\,\pi_j(L_\bullet)|=\mathrm{rg}\,\pi_j(L_\bullet)}$} \deleted{$\unsure{(j\geq0)}$} $(j\geq0)$ et une fonction \inserted{polynômiale} $f(n)$ à valeurs entières définie \inserted{sur les entiers} \deleted{pour $n$ grand}, à quelles conditions existe-t-il $L_\bullet$ très spécial avec $\nu_j=\mathrm{rg}\,\pi_j(L_\bullet)$ $(j\geq0)$ et $f(n)=f_{L_\bullet}(n)$ pour $n$ grand ? Notons que $g=f-\sum_{j\geq0}\nu_jP_j$ est encore une fonction \deleted{\unsure{entière}} \inserted{polynômiale} à valeurs entières, donc de la forme $\sum_{i\geq0}\mu_i\binom{n+1}{i}$ où $(\mu_i)\in\mathbb Z^{(\mathbb N)}$ est uniquement déterminé. Alors la condition demandée est
\[
\mu_0=0,\quad\mu_i\geq0\quad\text{si }i\geq1.
\]
Notons que \deleted{$\unsure{\mu_0=|\mu_0}$}
\[
\mu_0=g(-1)=f(-1)-\sum_{j\geq0}\nu_j\frac{(-1)(-2)\cdots(-j)}{j!}
\]
donc on a \unsure{bien la|encore}
\[
(\mu_0=0)\Longleftrightarrow\left(f(-1)=\underset{\substack{\Vert\\\chi((\nu_j))}}{\sum(-1)^j\nu_j}\right)
\]

\deleted{Déf. \deleted{Complexes} $k$-Modules $\frac12$ simpliciaux $L_\bullet$ \deleted{\unsure{prép.}} \inserted{\unsure{str.}} ``chaîne-parfait'' \inserted{(strictly chain-perfect)} $={}$ tels que $L_\bullet^!$ soit un complexe de chaînes de $k$-modules \deleted{\unsure{str.}} parfait (str. parfait) $={}$}

% archive sheet number: 12
